% ID: SEL-D0017-BC-E05
% ESTADO: cotejada. Segunda pasada visual completa y borrador compilado revisados.
% Fuente: Selectividad/matematicas_aplicadas/matematicas_aplicadas_P_E.pdf
% SHA-256: bc88bfeb0f184eb0c806c4efe04e84c24f3e5fcef9caaa12e94fd48d9f67a389
% Página: 2; bloque C; ejercicio original 5.

Sean $A$ y $B$ dos sucesos de un mismo experimento aleatorio.
\begin{enumerate}[label=\alph*)]
  \item \textbf{(0.5 puntos)} Si $P(A)\neq 0$ y $P(B)\neq 0$,
  \textquestiondown{}pueden ser los sucesos $A$ y $B$ independientes e
  incompatibles a la vez? Justifique la respuesta.
  \item \textbf{(2 puntos)} Sabiendo que $P(A)=0.3$, $P(B)=0.5$ y
  $P(A/B)=0.2$, calcule las siguientes probabilidades:
  \[
  P(A\cap B)\qquad P(A\cup B)\qquad P(A^c\cup B^c)\qquad P(A-B)
  \]
\end{enumerate}
